% =====================================================================
\begin{thebibliography}{99}
\addcontentsline{toc}{chapter}{Bibliografía}
\markboth{Bibliografía}{}
\small

\bibitem{ct}
T.~M. Cover y J.~A. Thomas.
\newblock \emph{Elements of Information Theory}, 2.ª ed.
\newblock Wiley, 2006.
\newblock Cap.~2, «Entropy, Relative Entropy, and Mutual Information».

\bibitem{fan}
A.~Fan, M.~Lewis e Y.~Dauphin.
\newblock Hierarchical neural story generation.
\newblock En \emph{Proceedings of the 56th Annual Meeting of the Association for Computational Linguistics (ACL)}, 2018.
\newblock \url{https://arxiv.org/abs/1805.04833}

\bibitem{gemini}
Gemini Team, Google.
\newblock Gemini: A family of highly capable multimodal models.
\newblock Informe técnico, 2023.
\newblock \url{https://arxiv.org/abs/2312.11805}

\bibitem{gbc}
I.~Goodfellow, Y.~Bengio y A.~Courville.
\newblock \emph{Deep Learning}.
\newblock MIT Press, 2016.
\newblock Cap.~7, «Regularization for Deep Learning»; cap.~8, «Optimization for Training Deep Models»; cap.~11, «Practical Methodology».
\newblock \url{https://www.deeplearningbook.org/}

\bibitem{holtzman}
A.~Holtzman, J.~Buys, L.~Du, M.~Forbes e Y.~Choi.
\newblock The curious case of neural text degeneration.
\newblock En \emph{International Conference on Learning Representations (ICLR)}, 2020.
\newblock \url{https://arxiv.org/abs/1904.09751}

\bibitem{nanogpt}
A.~Karpathy.
\newblock nanoGPT.
\newblock Repositorio de software, GitHub, 2022--2025. Archivos \texttt{model.py}, \texttt{train.py}, \texttt{sample.py} y \texttt{config/train\_shakespeare\_char.py}. Desde noviembre de 2025 su autor lo declara obsoleto y remite a nanochat.
\newblock \url{https://github.com/karpathy/nanoGPT}

\bibitem{llama3}
Llama Team, AI @ Meta.
\newblock The Llama 3 herd of models.
\newblock Informe técnico, 2024.
\newblock \url{https://arxiv.org/abs/2407.21783}

\bibitem{sgdr}
I.~Loshchilov y F.~Hutter.
\newblock SGDR: Stochastic gradient descent with warm restarts.
\newblock En \emph{International Conference on Learning Representations (ICLR)}, 2017.
\newblock \url{https://arxiv.org/abs/1608.03983}

\bibitem{adamw}
I.~Loshchilov y F.~Hutter.
\newblock Decoupled weight decay regularization.
\newblock En \emph{International Conference on Learning Representations (ICLR)}, 2019.
\newblock \url{https://arxiv.org/abs/1711.05101}

\bibitem{gpt2}
A.~Radford, J.~Wu, R.~Child, D.~Luan, D.~Amodei e I.~Sutskever.
\newblock Language models are unsupervised multitask learners.
\newblock Informe técnico, OpenAI, 2019.

\bibitem{cs336}
Stanford CS336.
\newblock Language Modeling from Scratch.
\newblock Stanford University, primavera de 2026. Tarea~1 (\emph{basics}), «Building a Transformer LM», versión~26.0.3.
\newblock \url{https://github.com/stanford-cs336/assignment1-basics}

\bibitem{llama}
H.~Touvron \emph{et al.}
\newblock LLaMA: Open and efficient foundation language models.
\newblock Prepublicación, 2023.
\newblock \url{https://arxiv.org/abs/2302.13971}

\bibitem{vaswani}
A.~Vaswani, N.~Shazeer, N.~Parmar, J.~Uszkoreit, L.~Jones, A.~N. Gomez, Ł.~Kaiser e I.~Polosukhin.
\newblock Attention is all you need.
\newblock En \emph{Advances in Neural Information Processing Systems 30 (NIPS)}, 2017.
\newblock \url{https://arxiv.org/abs/1706.03762}

\bibitem{grok1}
xAI.
\newblock Grok-1: pesos abiertos y código de ejemplo.
\newblock Repositorio de software, GitHub, 2024.
\newblock \url{https://github.com/xai-org/grok-1}

\bibitem{cs224n}
D.~Yang y Y.~Choi.
\newblock CS224N: Natural Language Processing with Deep Learning.
\newblock Stanford University, invierno de 2026. Clase~5, «Transformers»; clase~7, «Pretraining (Scaling, Systems, Data)».
\newblock \url{https://web.stanford.edu/class/cs224n/}

\bibitem{d2l}
A.~Zhang, Z.~C. Lipton, M.~Li y A.~J. Smola.
\newblock \emph{Dive into Deep Learning}, versión 1.0.3.
\newblock Cambridge University Press, 2023.
\newblock Cap.~9, «Recurrent Neural Networks» (§9.3 y §9.5); cap.~11, «Attention Mechanisms and Transformers» (§11.9).
\newblock \url{https://d2l.ai/}

\end{thebibliography}
